\chapter{矩阵求导}
\label{ch:matderiv}

% ════════════════════════════════════════════════════════════════
%  附录：矩阵求导（从零、自包含）
%  · 全书唯一的矩阵微积分教学单元，作为 ch:rigid_body / ch:lie /
%    ch:slam_est / ch:nlopt 中雅可比、最小二乘、扰动求导的统一地基。
%  · 本书统一取「分子布局」，与 ch:lie 的右扰动雅可比排布一致。
%  · 新 label 一律 app:matderiv-* / eq:matderiv-*；\cref 仅指向确已存在的章标签。
% ════════════════════════════════════════════════════════════════

本书从头到尾在做同一件事：把带噪声的传感器数据拟合到一个几何/概率模型上，再求解。无论是最小二乘的正规方程（\cref{ch:nlopt}）、卡尔曼滤波的协方差传播（\cref{ch:slam_est}），还是李群上的扰动求导（\cref{ch:lie}），核心动作都是\emph{对一个向量值或标量值函数关于一组参数求导}，得到一个被称为\textbf{雅可比矩阵}的对象。前面各章在用到时都默认读者已掌握这套运算；本附录把它从零补齐，给出统一的定义、排布约定与一张可直接套用的公式表，使全书在这一点上真正自包含。

读这一章不需要任何前置，只用到一元微积分里的偏导数 $\partial f/\partial x_i$ 和线性代数里的矩阵乘法、转置。读完后，\cref{ch:rigid_body} 的坐标变换求导、\cref{ch:lie} 的右扰动雅可比、\cref{ch:slam_est} 与 \cref{ch:nlopt} 的最小二乘与白化，都能落到同一套记号上。

% ════════════════════════════════════════════════════════════════
\section{动机：雅可比无处不在}
\label{sec:matderiv-motiv}

\paragraph{为什么估计问题离不开求导。}
机器人状态估计的主线是\emph{最小二乘}：给定残差 $\mathbf{e}(\mathbf{x})\in\mathbb{R}^n$（观测减预测），找参数 $\mathbf{x}\in\mathbb{R}^m$ 使 $J(\mathbf{x})=\tfrac12\|\mathbf{e}(\mathbf{x})\|^2$ 最小（\cref{ch:nlopt}）。求极小的第一步是令梯度为零；而残差通常非线性（相机投影、李群作用、$\cos/\operatorname{atan2}$），于是要在工作点把它\emph{线性化}：
\begin{equation}
  \mathbf{e}(\mathbf{x}+\Delta\mathbf{x})\approx\mathbf{e}(\mathbf{x})+\mathbf{J}\,\Delta\mathbf{x},
  \qquad
  \mathbf{J}=\frac{\partial\mathbf{e}}{\partial\mathbf{x}}\in\mathbb{R}^{n\times m}.
  \label{eq:matderiv-linearize}
\end{equation}
这里的 $\mathbf{J}$ 就是\emph{向量函数对向量求导}得到的雅可比矩阵——它是高斯牛顿正规方程（\cref{ch:nlopt} 的 \cref{eq:gn-normal}）的唯一原料。同样地：

\begin{itemize}
  \item \textbf{卡尔曼/扩展卡尔曼滤波}：非线性运动/观测模型 $f,h$ 的线性化用其雅可比 $\mathbf{F}=\partial f/\partial\mathbf{x}$、$\mathbf{H}=\partial h/\partial\mathbf{x}$；协方差经 $\mathbf{F}\boldsymbol{\Sigma}\mathbf{F}^\top$ 传播（\cref{ch:slam_est}）。
  \item \textbf{李群上的优化}：旋转作用点 $\mathbf{R}\mathbf{p}$ 关于扰动 $\boldsymbol{\varphi}$ 的导数 $\partial(\mathbf{R}\mathbf{p})/\partial\boldsymbol{\varphi}=-(\mathbf{R}\mathbf{p})^\wedge$（\cref{ch:lie}），本质上也是“向量函数对向量求导”，只是自变量住在切空间里。
  \item \textbf{相机标定、PnP、ICP}：每一个都要对残差求雅可比再解最小二乘（\cref{ch:rigid_body}、\cref{ch:nlopt}）。
\end{itemize}

\paragraph{反面：不统一约定会怎样。}
雅可比是个矩阵，它的“行对应谁、列对应谁”有两种摆法（分子布局 vs 分母布局，见 \cref{sec:matderiv-layout}）。两种摆法互为转置，单独看都自洽，\emph{但混用必出错}：海森近似 $\mathbf{J}^\top\mathbf{J}$ 会变成 $\mathbf{J}\mathbf{J}^\top$、维度对不上、正规方程左右不匹配。这类“转置 bug”在 SLAM 代码里极常见且难定位。因此本附录第一要务不是记住公式，而是\textbf{钉死一套排布约定并贯穿全书}。

\begin{insight}
一句话主线：\textbf{建模 $\to$ 写残差 $\to$ 对参数求雅可比 $\to$ 组装并解最小二乘/滤波更新}。整条链路上唯一“新”的数学动作就是\emph{求雅可比}。把它的定义和排布弄清楚，后面各章的求导都只是“代入公式”。
\end{insight}

% ════════════════════════════════════════════════════════════════
\section{标量对向量求导：梯度}
\label{sec:matderiv-scalar}

先从最简单的情形入手：自变量是向量、函数取标量值。设 $\mathbf{x}=[x_1,\dots,x_m]^\top\in\mathbb{R}^m$（本书约定向量为\emph{列}向量），$f(\mathbf{x})\in\mathbb{R}$ 为标量函数。把 $f$ 对每个分量 $x_i$ 求偏导，再把这 $m$ 个偏导排成一个向量，就得到\emph{梯度}。排成列还是排成行，是约定问题——两种写法都要给出，因为不同文献都在用。

\begin{definition}[标量对向量求导（梯度）]\label{def:matderiv-grad}
设 $f:\mathbb{R}^m\to\mathbb{R}$ 可微，$\mathbf{x}=[x_1,\dots,x_m]^\top$。
\textbf{对列向量 $\mathbf{x}$ 求导}得到一个\emph{列}向量（即通常意义的梯度 $\nabla f$）：
\begin{equation}
  \frac{\partial f}{\partial\mathbf{x}}
  =\Big[\tfrac{\partial f}{\partial x_1},\ \dots,\ \tfrac{\partial f}{\partial x_m}\Big]^\top
  =\nabla f\in\mathbb{R}^{m\times 1}.
  \label{eq:matderiv-grad-col}
\end{equation}
\textbf{对行向量 $\mathbf{x}^\top$ 求导}得到一个\emph{行}向量，恰为前者的转置：
\begin{equation}
  \frac{\partial f}{\partial\mathbf{x}^\top}
  =\Big[\tfrac{\partial f}{\partial x_1},\ \dots,\ \tfrac{\partial f}{\partial x_m}\Big]
  =\Big(\frac{\partial f}{\partial\mathbf{x}}\Big)^{\!\top}\in\mathbb{R}^{1\times m}.
  \label{eq:matderiv-grad-row}
\end{equation}
\end{definition}

\begin{note}
\textbf{记法与术语提醒.}\;名字在不同领域并不统一：优化与机器学习多把\emph{列}向量 $\partial f/\partial\mathbf{x}$ 叫\emph{梯度} $\nabla f$；而把\emph{行}向量 $\partial f/\partial\mathbf{x}^\top$ 看作 $1\times m$ 雅可比。两者只差一个转置，但代入二次型、链式法则时维度敏感，务必先认清当前用的是哪一种。本书在 \cref{ch:nlopt} 里目标函数的梯度记 $\nabla J$、海森记 $\mathbf{H}_J$，都按 \cref{def:matderiv-grad} 的\emph{列}向量约定。
\end{note}

\paragraph{几何意义。}
梯度 $\nabla f$ 指向 $f$ 上升最快的方向，其负方向 $-\nabla f$ 即最速下降方向——这正是 \cref{ch:nlopt} 梯度下降法的依据。在极小点处 $\nabla f=\mathbf{0}$，这条“一阶必要条件”是后面推正规方程的出发点。

% ════════════════════════════════════════════════════════════════
\section{向量对向量求导：雅可比矩阵}
\label{sec:matderiv-jacobian}

现在让函数也取向量值。设 $\mathbf{F}:\mathbb{R}^m\to\mathbb{R}^n$，写成分量
\begin{equation}
  \mathbf{F}(\mathbf{x})=\big[f_1(\mathbf{x}),\ \dots,\ f_n(\mathbf{x})\big]^\top,
  \label{eq:matderiv-vecfun}
\end{equation}
其中每个 $f_k:\mathbb{R}^m\to\mathbb{R}$ 是标量分量函数。把它对 $\mathbf{x}$ 求导，得到的对象由 $n\times m$ 个偏导 $\partial f_k/\partial x_l$ 组成，排成矩阵即雅可比。

\begin{definition}[向量对向量求导（雅可比矩阵）]\label{def:matderiv-jacobian}
设 $\mathbf{F}=[f_1,\dots,f_n]^\top:\mathbb{R}^m\to\mathbb{R}^n$ 可微，$\mathbf{x}=[x_1,\dots,x_m]^\top$。\textbf{把列函数 $\mathbf{F}$ 对行向量 $\mathbf{x}^\top$ 求导}，定义其\emph{雅可比矩阵}为
\begin{equation}
  \frac{\partial\mathbf{F}}{\partial\mathbf{x}^\top}
  =\begin{bmatrix}
    \dfrac{\partial f_1}{\partial\mathbf{x}^\top}\\[1.2ex]
    \vdots\\[0.6ex]
    \dfrac{\partial f_n}{\partial\mathbf{x}^\top}
  \end{bmatrix}
  =\begin{bmatrix}
    \dfrac{\partial f_1}{\partial x_1} & \dfrac{\partial f_1}{\partial x_2} & \cdots & \dfrac{\partial f_1}{\partial x_m}\\[1.4ex]
    \dfrac{\partial f_2}{\partial x_1} & \dfrac{\partial f_2}{\partial x_2} & \cdots & \dfrac{\partial f_2}{\partial x_m}\\[1.4ex]
    \vdots & \vdots & \ddots & \vdots\\[0.6ex]
    \dfrac{\partial f_n}{\partial x_1} & \dfrac{\partial f_n}{\partial x_2} & \cdots & \dfrac{\partial f_n}{\partial x_m}
  \end{bmatrix}\in\mathbb{R}^{n\times m}.
  \label{eq:matderiv-jacobian}
\end{equation}
其第 $(k,l)$ 元为 $\big[\partial\mathbf{F}/\partial\mathbf{x}^\top\big]_{kl}=\partial f_k/\partial x_l$：\textbf{行号 $k$ 随分子（函数分量），列号 $l$ 随分母（自变量分量）}。它把 \cref{eq:matderiv-grad-row} 的标量情形（$n=1$）推广到向量值函数。
\end{definition}

\paragraph{为什么是这个形状。}
\cref{eq:matderiv-jacobian} 让雅可比正好充当“线性化的乘子矩阵”：一阶泰勒展开
\begin{equation}
  \mathbf{F}(\mathbf{x}+\Delta\mathbf{x})\approx\mathbf{F}(\mathbf{x})+\frac{\partial\mathbf{F}}{\partial\mathbf{x}^\top}\,\Delta\mathbf{x},
  \label{eq:matderiv-taylor}
\end{equation}
左右维度自洽：$\mathbf{F}\in\mathbb{R}^n$、$\Delta\mathbf{x}\in\mathbb{R}^m$、雅可比 $\in\mathbb{R}^{n\times m}$，相乘得 $\mathbb{R}^n$。这正是 \cref{eq:matderiv-linearize} 残差线性化所用的形状，也是“雅可比 $=$ 局部最佳线性近似”的精确含义。每一行是一个标量分量 $f_k$ 的梯度（行向量），每一列描述所有分量对同一个 $x_l$ 的敏感度。

\begin{note}
\textbf{先看一个例子.}\;取 $\mathbf{F}(\mathbf{x})=\mathbf{A}\mathbf{x}$，$\mathbf{A}\in\mathbb{R}^{n\times m}$ 为常矩阵。第 $k$ 个分量 $f_k=\sum_{l}A_{kl}x_l$，故 $\partial f_k/\partial x_l=A_{kl}$，按 \cref{def:matderiv-jacobian} 排出来就是 $\mathbf{A}$ 本身。这给出下一节要正式列出的规范例 $\partial(\mathbf{A}\mathbf{x})/\partial\mathbf{x}^\top=\mathbf{A}$。
\end{note}

% ════════════════════════════════════════════════════════════════
\section{分子布局与分母布局}
\label{sec:matderiv-layout}

\cref{def:matderiv-jacobian} 把行系于分子、列系于分母——这只是\emph{两种}主流约定之一。另一种把行列对调。两者互为转置，使用时必须二选一并贯穿到底，否则维度与转置全乱。这是矩阵求导里最大的“坑”，单列一节讲清。

\begin{definition}[分子布局与分母布局]\label{def:matderiv-layout}
对向量函数 $\mathbf{F}\in\mathbb{R}^n$ 关于向量 $\mathbf{x}\in\mathbb{R}^m$ 求导，导数矩阵有两种排布约定\cite{petersen2012cookbook,magnus2019matrix}：
\begin{itemize}
  \item \textbf{分子布局（numerator layout，又称 Jacobian layout）}：结果形状\emph{随分子}，为 $n\times m$；第 $(k,l)$ 元 $=\partial F_k/\partial x_l$。即 \cref{eq:matderiv-jacobian}，记作 $\dfrac{\partial\mathbf{F}}{\partial\mathbf{x}^\top}$。
  \item \textbf{分母布局（denominator layout，又称 gradient layout / Hessian formulation）}：结果形状\emph{随分母}，为 $m\times n$；第 $(l,k)$ 元 $=\partial F_k/\partial x_l$。它恰为分子布局的转置，记作 $\dfrac{\partial\mathbf{F}}{\partial\mathbf{x}}=\Big(\dfrac{\partial\mathbf{F}}{\partial\mathbf{x}^\top}\Big)^{\!\top}$。
\end{itemize}
两者关系：
\begin{equation}
  \underbrace{\frac{\partial\mathbf{F}}{\partial\mathbf{x}}}_{\text{分母布局},\ m\times n}
  =\bigg(\underbrace{\frac{\partial\mathbf{F}}{\partial\mathbf{x}^\top}}_{\text{分子布局},\ n\times m}\bigg)^{\!\top}.
  \label{eq:matderiv-layout-rel}
\end{equation}
\end{definition}

\begin{table}[H]\centering\small
\caption{分子布局与分母布局对照（$\mathbf{F}\in\mathbb{R}^n$、$\mathbf{x}\in\mathbb{R}^m$、$f\in\mathbb{R}$ 标量）}
\label{app:matderiv-layout-table}
\begin{tabular}{lll}
\toprule
对象 & 分子布局（本书） & 分母布局 \\
\midrule
标量 $f$ 对向量 & $\dfrac{\partial f}{\partial\mathbf{x}^\top}\in\mathbb{R}^{1\times m}$（行向量） & $\dfrac{\partial f}{\partial\mathbf{x}}\in\mathbb{R}^{m\times 1}$（列向量，即 $\nabla f$）\\[1.6ex]
向量 $\mathbf{F}$ 对向量 & $\dfrac{\partial\mathbf{F}}{\partial\mathbf{x}^\top}\in\mathbb{R}^{n\times m}$（雅可比） & $\dfrac{\partial\mathbf{F}}{\partial\mathbf{x}}\in\mathbb{R}^{m\times n}$ \\[1.6ex]
泰勒展开乘子 & $\mathbf{F}(\mathbf{x}{+}\Delta\mathbf{x})\approx\mathbf{F}{+}\dfrac{\partial\mathbf{F}}{\partial\mathbf{x}^\top}\Delta\mathbf{x}$ & $\mathbf{F}{+}\Big(\dfrac{\partial\mathbf{F}}{\partial\mathbf{x}}\Big)^{\!\top}\Delta\mathbf{x}$ \\[1.6ex]
链式法则方向 & 由外到内\emph{左乘}（\cref{thm:matderiv-chain}） & 由内到外\emph{右乘} \\
\bottomrule
\end{tabular}
\end{table}

\begin{insight}[本书的统一约定]
\textbf{本书对“向量对向量”求导一律取分子布局}：向量函数对向量求导得 $n\times m$ 的雅可比 $\partial\mathbf{F}/\partial\mathbf{x}^\top$（行随分子、列随分母）。（标量梯度 $\nabla f=\partial f/\partial\mathbf{x}$ 仍按惯例记为\emph{列}向量，即 $1\times m$ 分子布局行的转置——这是文献通行的混合约定，不影响下面雅可比的排布。）这与 \cref{ch:lie} 中“分子 $[\mathbf{a};\mathbf{b}]$ 对分母 $[\mathbf{x};\mathbf{y}]$ 求导写成分块矩阵、行随分子列随分母”的声明完全一致，也与十四讲\cite{gaoxiang2019slam14}、多数 SLAM 求解器（Ceres、g2o、GTSAM 把残差雅可比按 $\partial\mathbf{e}/\partial\mathbf{x}$、行$=$残差维、列$=$状态维组织）一致。本附录是 \cref{ch:lie} 那条声明的\emph{依据}：从这里起，全书的雅可比排布都按分子布局核对。
\end{insight}

\paragraph{“忽略分母转置”的简记由来。}
严格写分子布局时，向量对向量求导要在分母加转置号 $\partial\mathbf{F}/\partial\mathbf{x}^\top$，每次都写很繁琐。文献里有一个广泛采用的简记\cite{gaoxiang2019slam14}：在\emph{不引起歧义}时省略分母上的转置，把
\begin{equation}
  \frac{\partial\mathbf{F}}{\partial\mathbf{x}}\ \xleftarrow{\ \text{简记}\ }\ \frac{\partial\mathbf{F}}{\partial\mathbf{x}^\top}
  \quad(\text{当语境约定分子布局时}),
  \label{eq:matderiv-shorthand}
\end{equation}
即\emph{看到“向量对向量”的记号就默认它是 $n\times m$ 的雅可比}。本书在后续各章（如 \cref{ch:lie} 的 $\partial(\mathbf{R}\mathbf{p})/\partial\boldsymbol{\varphi}$、\cref{ch:nlopt} 的 $\partial\mathbf{e}/\partial\mathbf{x}$）正是这样用的：写 $\partial\mathbf{F}/\partial\mathbf{x}$，意思是分子布局的雅可比。

\begin{pitfall}[简记的陷阱：别把“忽略转置”当成“可以乱转置”]
简记 \cref{eq:matderiv-shorthand} \emph{只}是省去分母的书写转置号，\textbf{不改变矩阵的形状与含义}——它仍是 $n\times m$ 的分子布局雅可比。三个常见错误：
\begin{enumerate}
  \item \textbf{把简记当分母布局}：以为 $\partial\mathbf{F}/\partial\mathbf{x}$（无转置号）就是 $m\times n$，于是组装出 $\mathbf{J}\mathbf{J}^\top$ 而非 $\mathbf{J}^\top\mathbf{J}$，正规方程左右不匹配。
  \item \textbf{标量退化时摆错}：标量 $f$ 的“梯度”到底取列（$\partial f/\partial\mathbf{x}$）还是行（$\partial f/\partial\mathbf{x}^\top$），与向量情形的简记不是一回事；\cref{ch:nlopt} 统一把 $\nabla J$ 取\emph{列}。
  \item \textbf{跨书直接抄公式}：Barfoot、十四讲、Magnus--Neudecker 的转置位置可能不同。引用异制公式前，先判明它是分子还是分母布局，必要时按 \cref{eq:matderiv-layout-rel} 转一次再用。
\end{enumerate}
铁律与 \cref{ch:lie}/\cref{ch:nlopt} 的雅可比纪律相同：\textbf{全书只用一种布局（分子），混用必出维度/转置 bug}。
\end{pitfall}

% ════════════════════════════════════════════════════════════════
\section{常用求导公式表}
\label{sec:matderiv-formulas}

下面在分子布局下，给出 SLAM/估计里反复出现的几条恒等式，每条都附简短推导，确保可独立复现，而非死记。除特别说明外，$\mathbf{x}\in\mathbb{R}^m$ 为变量，$\mathbf{A},\mathbf{a}$ 为常量（与 $\mathbf{x}$ 无关），$f$ 为标量函数，$\mathbf{F}$ 为向量函数。

\subsection{线性型与规范例}
\label{sec:matderiv-linear}

\begin{proposition}[线性型的雅可比（规范例）]\label{prop:matderiv-Ax}
设 $\mathbf{A}\in\mathbb{R}^{n\times m}$ 为常矩阵，则
\begin{equation}
  \frac{\partial(\mathbf{A}\mathbf{x})}{\partial\mathbf{x}^\top}=\mathbf{A}\in\mathbb{R}^{n\times m}.
  \label{eq:matderiv-Ax}
\end{equation}
这是分子布局\emph{规范}的典型例：列函数 $\mathbf{A}\mathbf{x}$ 对行向量 $\mathbf{x}^\top$ 求导，结果恰为 $\mathbf{A}$。
\end{proposition}
\begin{proof}
第 $k$ 个分量 $(\mathbf{A}\mathbf{x})_k=\sum_{l=1}^{m}A_{kl}x_l$，对 $x_l$ 求偏导得常数 $A_{kl}$。按 \cref{def:matderiv-jacobian}，雅可比第 $(k,l)$ 元 $=\partial(\mathbf{A}\mathbf{x})_k/\partial x_l=A_{kl}$，即整块等于 $\mathbf{A}$。\hfill$\square$
\end{proof}

\begin{note}
\textbf{简记形式（务必给出）.}\;按 \cref{eq:matderiv-shorthand} 省略分母转置号，\cref{eq:matderiv-Ax} 常写作
\begin{equation}
  \frac{\partial(\mathbf{A}\mathbf{x})}{\partial\mathbf{x}}=\mathbf{A}.
  \label{eq:matderiv-Ax-short}
\end{equation}
两式\emph{同义}：等号右边都是 $n\times m$ 的 $\mathbf{A}$，只是左边是否写出分母转置号之别。本书在 \cref{ch:slam_est}/\cref{ch:nlopt} 默认用简记形式 \cref{eq:matderiv-Ax-short}，但其精确含义始终是 \cref{eq:matderiv-Ax}。
\end{note}

\begin{proposition}[行向量函数对列向量求导 $=$ 转置规则]\label{prop:matderiv-transpose}
设 $\mathbf{F}:\mathbb{R}^m\to\mathbb{R}^n$ 为列函数，则其转置 $\mathbf{F}^\top$（行向量函数）对列向量 $\mathbf{x}$ 求导，等于雅可比的转置：
\begin{equation}
  \frac{\partial\mathbf{F}^\top}{\partial\mathbf{x}}
  =\Big(\frac{\partial\mathbf{F}}{\partial\mathbf{x}^\top}\Big)^{\!\top}\in\mathbb{R}^{m\times n}.
  \label{eq:matderiv-transpose}
\end{equation}
特别地，由 \cref{prop:matderiv-Ax} 立得 $\dfrac{\partial(\mathbf{A}\mathbf{x})^\top}{\partial\mathbf{x}}=\dfrac{\partial(\mathbf{x}^\top\mathbf{A}^\top)}{\partial\mathbf{x}}=\mathbf{A}^\top$。
\end{proposition}
\begin{proof}
$\mathbf{F}^\top$ 是行向量，其分量与 $\mathbf{F}$ 相同；对列向量 $\mathbf{x}$ 求导时，按分子布局“行随分子分量、列随分母分量”，而此处分子已转成行、分母为列，于是行列与 \cref{eq:matderiv-jacobian} 恰好对调，得 $\big(\partial\mathbf{F}/\partial\mathbf{x}^\top\big)^\top$。这与布局关系 \cref{eq:matderiv-layout-rel} 一致：转置分子等价于把分母从行换成列。\hfill$\square$
\end{proof}

\subsection{二次型}
\label{sec:matderiv-quadratic}

二次型是最小二乘目标函数的“原子”——$\|\mathbf{e}\|^2=\mathbf{e}^\top\mathbf{e}$、马氏距离 $\mathbf{e}^\top\mathbf{W}^{-1}\mathbf{e}$ 都是二次型。下面三条由浅入深。

\begin{proposition}[内积型 $\partial(\mathbf{x}^\top\mathbf{a})/\partial\mathbf{x}$]\label{prop:matderiv-xta}
设 $\mathbf{a}\in\mathbb{R}^m$ 为常向量，则标量 $\mathbf{x}^\top\mathbf{a}=\mathbf{a}^\top\mathbf{x}=\sum_l a_l x_l$ 满足
\begin{equation}
  \frac{\partial(\mathbf{x}^\top\mathbf{a})}{\partial\mathbf{x}}=\frac{\partial(\mathbf{a}^\top\mathbf{x})}{\partial\mathbf{x}}=\mathbf{a},
  \qquad
  \frac{\partial(\mathbf{x}^\top\mathbf{a})}{\partial\mathbf{x}^\top}=\mathbf{a}^\top.
  \label{eq:matderiv-xta}
\end{equation}
\end{proposition}
\begin{proof}
$\partial(\sum_l a_l x_l)/\partial x_i=a_i$。按 \cref{def:matderiv-grad}，对列向量 $\mathbf{x}$ 求导得列向量 $\mathbf{a}$；对行向量 $\mathbf{x}^\top$ 求导得行向量 $\mathbf{a}^\top$。\hfill$\square$
\end{proof}

\begin{proposition}[平方型 $\partial(\mathbf{x}^\top\mathbf{x})/\partial\mathbf{x}$]\label{prop:matderiv-xtx}
$\mathbf{x}^\top\mathbf{x}=\|\mathbf{x}\|^2=\sum_l x_l^2$ 满足
\begin{equation}
  \frac{\partial(\mathbf{x}^\top\mathbf{x})}{\partial\mathbf{x}}=2\mathbf{x}.
  \label{eq:matderiv-xtx}
\end{equation}
\end{proposition}
\begin{proof}
$\partial(\sum_l x_l^2)/\partial x_i=2x_i$，排成列向量即 $2\mathbf{x}$。这是 \cref{prop:matderiv-xtAx} 取 $\mathbf{A}=\mathbf{I}$ 的特例（$(\mathbf{I}+\mathbf{I}^\top)\mathbf{x}=2\mathbf{x}$）。\hfill$\square$
\end{proof}

\begin{proposition}[一般二次型 $\partial(\mathbf{x}^\top\mathbf{A}\mathbf{x})/\partial\mathbf{x}$]\label{prop:matderiv-xtAx}
设 $\mathbf{A}\in\mathbb{R}^{m\times m}$ 为常矩阵，则标量 $\mathbf{x}^\top\mathbf{A}\mathbf{x}$ 满足\cite{petersen2012cookbook}
\begin{equation}
  \frac{\partial(\mathbf{x}^\top\mathbf{A}\mathbf{x})}{\partial\mathbf{x}}=(\mathbf{A}+\mathbf{A}^\top)\mathbf{x}.
  \label{eq:matderiv-xtAx}
\end{equation}
特别地，$\mathbf{A}$ \emph{对称}（$\mathbf{A}=\mathbf{A}^\top$，如协方差/信息矩阵）时 $\dfrac{\partial(\mathbf{x}^\top\mathbf{A}\mathbf{x})}{\partial\mathbf{x}}=2\mathbf{A}\mathbf{x}$。
\end{proposition}
\begin{proof}
按分量展开 $\mathbf{x}^\top\mathbf{A}\mathbf{x}=\sum_{p}\sum_{q}A_{pq}x_p x_q$。对 $x_i$ 求偏导，只有含 $x_i$ 的项存活，分两类：当 $p=i$（$q$ 任意）贡献 $\sum_q A_{iq}x_q$；当 $q=i$（$p$ 任意）贡献 $\sum_p A_{pi}x_p$（$p=q=i$ 的项 $A_{ii}x_i^2$ 求导得 $2A_{ii}x_i$，正好被两类各计一次）。合并
\begin{equation*}
  \frac{\partial(\mathbf{x}^\top\mathbf{A}\mathbf{x})}{\partial x_i}
  =\sum_q A_{iq}x_q+\sum_p A_{pi}x_p
  =(\mathbf{A}\mathbf{x})_i+(\mathbf{A}^\top\mathbf{x})_i
  =\big((\mathbf{A}+\mathbf{A}^\top)\mathbf{x}\big)_i .
\end{equation*}
排成列向量即 \cref{eq:matderiv-xtAx}。\hfill$\square$
\end{proof}

\begin{pitfall}[二次型求导漏掉 $\mathbf{A}^\top$]
最常见的错误是把 $\partial(\mathbf{x}^\top\mathbf{A}\mathbf{x})/\partial\mathbf{x}$ 直接写成 $\mathbf{A}\mathbf{x}$（少了一半）。\emph{只有} $\mathbf{A}$ 对称时才有 $2\mathbf{A}\mathbf{x}$；$\mathbf{A}$ 非对称时必须是 $(\mathbf{A}+\mathbf{A}^\top)\mathbf{x}$。所幸 SLAM 里二次型几乎总以信息矩阵 $\boldsymbol{\Sigma}^{-1}$（对称正定）加权出现，故常见的 $\partial(\mathbf{e}^\top\boldsymbol{\Sigma}^{-1}\mathbf{e})/\partial\mathbf{e}=2\boldsymbol{\Sigma}^{-1}\mathbf{e}$ 是对的——但这是“对称”帮了忙，不是普遍规律。
\end{pitfall}

\subsection{链式法则与复合}
\label{sec:matderiv-chain}

实际残差几乎都是复合函数（先变换、再投影、再相减）。链式法则是把它们拆开逐段求导的工具，也是 \cref{ch:nlopt} 组装雅可比的核心。

\begin{theorem}[向量复合的链式法则]\label{thm:matderiv-chain}
设 $\mathbf{u}=\mathbf{g}(\mathbf{x}):\mathbb{R}^m\to\mathbb{R}^p$，$\mathbf{F}=\mathbf{h}(\mathbf{u}):\mathbb{R}^p\to\mathbb{R}^n$，复合 $\mathbf{F}(\mathbf{x})=\mathbf{h}(\mathbf{g}(\mathbf{x}))$。在分子布局下，复合的雅可比等于各段雅可比\emph{从外到内左乘}：
\begin{equation}
  \frac{\partial\mathbf{F}}{\partial\mathbf{x}^\top}
  =\frac{\partial\mathbf{h}}{\partial\mathbf{u}^\top}\,\frac{\partial\mathbf{g}}{\partial\mathbf{x}^\top}
  \in\mathbb{R}^{n\times m},
  \label{eq:matderiv-chain}
\end{equation}
其中 $\partial\mathbf{h}/\partial\mathbf{u}^\top\in\mathbb{R}^{n\times p}$、$\partial\mathbf{g}/\partial\mathbf{x}^\top\in\mathbb{R}^{p\times m}$，相乘维度 $n\times m$ 自洽。
\end{theorem}
\begin{proof}
取标量分量 $F_k=h_k(\mathbf{u})$，$u_j=g_j(\mathbf{x})$。一元/多元链式法则给
\begin{equation*}
  \frac{\partial F_k}{\partial x_l}=\sum_{j=1}^{p}\frac{\partial h_k}{\partial u_j}\frac{\partial g_j}{\partial x_l}.
\end{equation*}
右端正是矩阵乘积 $\big(\partial\mathbf{h}/\partial\mathbf{u}^\top\big)\big(\partial\mathbf{g}/\partial\mathbf{x}^\top\big)$ 的第 $(k,l)$ 元。逐元相等即 \cref{eq:matderiv-chain}。\hfill$\square$
\end{proof}

\begin{note}
\textbf{布局决定乘法方向.}\;分子布局下链式法则“由外到内\emph{左}乘”（外层雅可比在左），与一元链式法则 $\tfrac{\mathrm{d}F}{\mathrm{d}x}=\tfrac{\mathrm{d}F}{\mathrm{d}u}\tfrac{\mathrm{d}u}{\mathrm{d}x}$ 的书写顺序一致，自然好记；若用分母布局则顺序相反（由内到外右乘）。这又是“坚持一种布局”的好处之一。
\end{note}

\begin{corollary}[复合线性映射 $\partial(\mathbf{A}\mathbf{u})/\partial\mathbf{x}$]\label{cor:matderiv-Au}
设 $\mathbf{u}=\mathbf{u}(\mathbf{x}):\mathbb{R}^m\to\mathbb{R}^p$ 可微，$\mathbf{A}\in\mathbb{R}^{n\times p}$ 常矩阵，则
\begin{equation}
  \frac{\partial(\mathbf{A}\mathbf{u})}{\partial\mathbf{x}^\top}
  =\mathbf{A}\,\frac{\partial\mathbf{u}}{\partial\mathbf{x}^\top}\in\mathbb{R}^{n\times m}.
  \label{eq:matderiv-Au}
\end{equation}
\end{corollary}
\begin{proof}
在 \cref{thm:matderiv-chain} 中取外层 $\mathbf{h}(\mathbf{u})=\mathbf{A}\mathbf{u}$，由 \cref{prop:matderiv-Ax} 知 $\partial\mathbf{h}/\partial\mathbf{u}^\top=\mathbf{A}$，代入即得。这条在 \cref{ch:slam_est} 协方差传播、\cref{ch:nlopt} 残差雅可比组装里反复出现（如 $\mathbf{e}=\mathbf{z}-\mathbf{A}\mathbf{u}(\mathbf{x})$ 的雅可比 $\partial\mathbf{e}/\partial\mathbf{x}^\top=-\mathbf{A}\,\partial\mathbf{u}/\partial\mathbf{x}^\top$）。\hfill$\square$
\end{proof}

\begin{table}[H]\centering\small
\caption{常用矩阵求导公式速查（分子布局；$\mathbf{A},\mathbf{a}$ 为常量，$\mathbf{x}$ 为变量）}
\label{app:matderiv-formula-table}
\begin{tabular}{lll}
\toprule
表达式 & 导数 & 出处 \\
\midrule
$\mathbf{A}\mathbf{x}$（向量） & $\dfrac{\partial(\mathbf{A}\mathbf{x})}{\partial\mathbf{x}^\top}=\mathbf{A}$ & \cref{prop:matderiv-Ax} \\[1.4ex]
$(\mathbf{A}\mathbf{x})^\top=\mathbf{x}^\top\mathbf{A}^\top$ & $\dfrac{\partial(\mathbf{x}^\top\mathbf{A}^\top)}{\partial\mathbf{x}}=\mathbf{A}^\top$ & \cref{prop:matderiv-transpose} \\[1.4ex]
$\mathbf{x}^\top\mathbf{a}=\mathbf{a}^\top\mathbf{x}$（标量） & $\dfrac{\partial}{\partial\mathbf{x}}=\mathbf{a}$ & \cref{prop:matderiv-xta} \\[1.4ex]
$\mathbf{x}^\top\mathbf{x}=\|\mathbf{x}\|^2$（标量） & $\dfrac{\partial}{\partial\mathbf{x}}=2\mathbf{x}$ & \cref{prop:matderiv-xtx} \\[1.4ex]
$\mathbf{x}^\top\mathbf{A}\mathbf{x}$（标量） & $\dfrac{\partial}{\partial\mathbf{x}}=(\mathbf{A}+\mathbf{A}^\top)\mathbf{x}$；对称时 $2\mathbf{A}\mathbf{x}$ & \cref{prop:matderiv-xtAx} \\[1.4ex]
$\mathbf{h}(\mathbf{g}(\mathbf{x}))$（复合） & $\dfrac{\partial\mathbf{h}}{\partial\mathbf{u}^\top}\dfrac{\partial\mathbf{g}}{\partial\mathbf{x}^\top}$（外到内左乘） & \cref{thm:matderiv-chain} \\[1.4ex]
$\mathbf{A}\mathbf{u}(\mathbf{x})$（复合线性） & $\mathbf{A}\dfrac{\partial\mathbf{u}}{\partial\mathbf{x}^\top}$ & \cref{cor:matderiv-Au} \\[1.4ex]
\midrule
\multicolumn{3}{l}{\emph{以下为选读（标量/矩阵对矩阵、代数补遗、离散化与统计变换）}}\\[0.4ex]
$\operatorname{tr}(\mathbf{A}\mathbf{X}),\ \ln\det\mathbf{X}$ & $\mathbf{A}^\top,\ \mathbf{X}^{-\top}$（迹技巧） & \cref{eq:matderiv-trace-formulas} \\[1.4ex]
$\operatorname{vec}(\mathbf{A}\mathbf{X}\mathbf{B})$ & $(\mathbf{B}^\top\otimes\mathbf{A})\operatorname{vec}(\mathbf{X})$，即 $\partial/\partial\operatorname{vec}(\mathbf{X})^\top=\mathbf{B}^\top\otimes\mathbf{A}$ & \cref{eq:matderiv-vec-AXB} \\[1.4ex]
对称 $\mathbf{X}$ 的 $\partial f/\partial\mathbf{X}$ & 取对称化导数 $\operatorname{sym}(\partial f/\partial\mathbf{X})$ & \cref{prop:matderiv-sym} \\[1.4ex]
分块求逆 / 舒尔补 & $\mathbf{M}^{-1}$ 经 $\mathbf{M}/\mathbf{D}=\mathbf{A}-\mathbf{B}\mathbf{D}^{-1}\mathbf{C}$（边缘化/SMW） & \cref{prop:matderiv-block-inv} \\[1.4ex]
LTI 离散化 $\boldsymbol{\Phi},\mathbf{Q}_k$ & Van Loan 增广阵一次指数读出 & \cref{eq:matderiv-vanloan-readout} \\[1.4ex]
FIM 重参数化 & $\boldsymbol{\mathcal{I}}_{\boldsymbol{\phi}}=\mathbf{J}^\top\boldsymbol{\mathcal{I}}_{\boldsymbol{\theta}}\mathbf{J}$（雅可比夹心） & \cref{thm:matderiv-fim-reparam} \\
\bottomrule
\end{tabular}
\end{table}

% ════════════════════════════════════════════════════════════════
\section{贯穿例：从二次型求导到正规方程}
\label{sec:matderiv-normal}

把上面的公式串起来用一次：推导线性最小二乘的\emph{正规方程}。这正是 \cref{ch:nlopt} 的 \cref{eq:nlopt-gls} 与 \cref{eq:gn-normal} 的源头，也是 \cref{ch:slam_est} 线性高斯估计的代数核心。

\begin{derivation}[最小二乘正规方程 $\mathbf{A}^\top\mathbf{A}\mathbf{x}=\mathbf{A}^\top\mathbf{b}$]
给定超定线性系统 $\mathbf{A}\mathbf{x}\approx\mathbf{b}$（$\mathbf{A}\in\mathbb{R}^{n\times m}$，$n>m$，列满秩），求使残差平方和最小的 $\mathbf{x}$：
\begin{equation}
  \min_{\mathbf{x}}\ J(\mathbf{x})=\|\mathbf{A}\mathbf{x}-\mathbf{b}\|^2
  =(\mathbf{A}\mathbf{x}-\mathbf{b})^\top(\mathbf{A}\mathbf{x}-\mathbf{b}).
  \label{eq:matderiv-ls-obj}
\end{equation}
\emph{Step 1（展开为二次型）.} 利用 $(\mathbf{A}\mathbf{x})^\top=\mathbf{x}^\top\mathbf{A}^\top$ 与转置分配律：
\begin{equation}
  J(\mathbf{x})
  =\mathbf{x}^\top\mathbf{A}^\top\mathbf{A}\mathbf{x}
  -\mathbf{x}^\top\mathbf{A}^\top\mathbf{b}
  -\mathbf{b}^\top\mathbf{A}\mathbf{x}
  +\mathbf{b}^\top\mathbf{b}.
  \label{eq:matderiv-ls-expand}
\end{equation}
注意中间两项互为转置且都是标量，故相等：$\mathbf{x}^\top\mathbf{A}^\top\mathbf{b}=\mathbf{b}^\top\mathbf{A}\mathbf{x}$，合并为 $-2\mathbf{b}^\top\mathbf{A}\mathbf{x}$。

\emph{Step 2（逐项求梯度，取列向量）.} 对 \cref{eq:matderiv-ls-expand} 逐项用公式表：
\begin{itemize}
  \item 二次项：$\mathbf{A}^\top\mathbf{A}$ 对称，由 \cref{prop:matderiv-xtAx} 得 $\dfrac{\partial}{\partial\mathbf{x}}\big(\mathbf{x}^\top\mathbf{A}^\top\mathbf{A}\mathbf{x}\big)=2\mathbf{A}^\top\mathbf{A}\mathbf{x}$；
  \item 线性项：把 $-2\mathbf{b}^\top\mathbf{A}\mathbf{x}=-2(\mathbf{A}^\top\mathbf{b})^\top\mathbf{x}$ 看作内积，由 \cref{prop:matderiv-xta} 得 $\dfrac{\partial}{\partial\mathbf{x}}=-2\mathbf{A}^\top\mathbf{b}$；
  \item 常数项 $\mathbf{b}^\top\mathbf{b}$ 梯度为零。
\end{itemize}
合得
\begin{equation}
  \frac{\partial J}{\partial\mathbf{x}}=2\mathbf{A}^\top\mathbf{A}\mathbf{x}-2\mathbf{A}^\top\mathbf{b}.
  \label{eq:matderiv-ls-grad}
\end{equation}

\emph{Step 3（一阶必要条件，令梯度为零）.} 凸二次目标的极小点处梯度为零：
\begin{equation}
  \frac{\partial J}{\partial\mathbf{x}}=\mathbf{0}
  \ \Longrightarrow\
  \boxed{\ \mathbf{A}^\top\mathbf{A}\,\mathbf{x}=\mathbf{A}^\top\mathbf{b}\ }
  \ \Longrightarrow\
  \mathbf{x}^*=(\mathbf{A}^\top\mathbf{A})^{-1}\mathbf{A}^\top\mathbf{b}
  \quad(\mathbf{A}\text{ 列满秩}).
  \label{eq:matderiv-normal}
\end{equation}
这条 $\mathbf{A}^\top\mathbf{A}\mathbf{x}=\mathbf{A}^\top\mathbf{b}$ 即\emph{正规方程}（normal equations）。$\mathbf{A}^\top\mathbf{A}$ 在 $\mathbf{A}$ 列满秩时对称正定可逆，故解唯一。加权（马氏）版只需把目标换成 $\|\mathbf{A}\mathbf{x}-\mathbf{b}\|_{\boldsymbol{\Sigma}^{-1}}^2$，同法得 $\mathbf{A}^\top\boldsymbol{\Sigma}^{-1}\mathbf{A}\mathbf{x}=\mathbf{A}^\top\boldsymbol{\Sigma}^{-1}\mathbf{b}$——正是 \cref{ch:nlopt} 的 \cref{eq:nlopt-gls}。
\end{derivation}

\begin{insight}
\cref{eq:matderiv-normal} 把本附录的零件全用上了：转置规则展开二次型（\cref{prop:matderiv-transpose}）、二次型求导（\cref{prop:matderiv-xtAx}）、内积求导（\cref{prop:matderiv-xta}）、令梯度为零（\cref{def:matderiv-grad}）。\cref{ch:nlopt} 的高斯牛顿正规方程 \cref{eq:gn-normal} 不过是把这里的常矩阵 $\mathbf{A}$ 换成残差雅可比 $\mathbf{J}$、在每个线性化点重复这一步。换言之，\textbf{整个最小二乘求解的代数骨架，就是“二次型对向量求导、置零”}。
\end{insight}

% ════════════════════════════════════════════════════════════════
\section{延伸：标量对矩阵、矩阵对矩阵（选读）}
\label{sec:matderiv-matrix}

前面的向量求导已足够支撑全书主线。但有些场合（如对协方差矩阵求导、对旋转矩阵元素求导）会遇到\emph{标量对矩阵}甚至\emph{矩阵对矩阵}求导。本节给出最小够用的结论与“迹技巧”，属选读，初读可跳过，用到时再回看。

\subsection{标量对矩阵求导}
\label{sec:matderiv-scalar-mat}

\begin{definition}[标量对矩阵求导]\label{def:matderiv-scalar-mat}
设 $f:\mathbb{R}^{p\times q}\to\mathbb{R}$ 为矩阵变量 $\mathbf{X}=[X_{ij}]$ 的标量函数。本书取\emph{分子布局}定义其导数为与 $\mathbf{X}$ \emph{同形}的矩阵，逐元为偏导：
\begin{equation}
  \Big[\frac{\partial f}{\partial\mathbf{X}}\Big]_{ij}=\frac{\partial f}{\partial X_{ij}},
  \qquad \frac{\partial f}{\partial\mathbf{X}}\in\mathbb{R}^{p\times q}.
  \label{eq:matderiv-scalar-mat}
\end{equation}
\end{definition}

\paragraph{迹技巧（trace trick）。}
对矩阵求导最顺手的工具是\emph{微分} $\mathrm{d}f$ 加\emph{迹}\cite{magnus2019matrix}。其要点（Magnus--Neudecker 的\emph{识别定理}思想）：若能把标量函数的微分整理成
\begin{equation}
  \mathrm{d}f=\operatorname{tr}\!\big(\mathbf{G}^\top\,\mathrm{d}\mathbf{X}\big),
  \quad\text{则}\quad
  \frac{\partial f}{\partial\mathbf{X}}=\mathbf{G}.
  \label{eq:matderiv-trace-id}
\end{equation}
配合迹的循环性 $\operatorname{tr}(\mathbf{A}\mathbf{B}\mathbf{C})=\operatorname{tr}(\mathbf{C}\mathbf{A}\mathbf{B})$ 与 $\operatorname{tr}(\mathbf{A})=\operatorname{tr}(\mathbf{A}^\top)$，可机械地导出常用结论：
\begin{equation}
  \frac{\partial\operatorname{tr}(\mathbf{A}\mathbf{X})}{\partial\mathbf{X}}=\mathbf{A}^\top,\qquad
  \frac{\partial\operatorname{tr}(\mathbf{X}^\top\mathbf{A}\mathbf{X})}{\partial\mathbf{X}}=(\mathbf{A}+\mathbf{A}^\top)\mathbf{X},\qquad
  \frac{\partial\ln\det\mathbf{X}}{\partial\mathbf{X}}=\mathbf{X}^{-\top}.
  \label{eq:matderiv-trace-formulas}
\end{equation}
\cref{eq:matderiv-trace-formulas} 第二式是 \cref{prop:matderiv-xtAx} 的矩阵推广（取 $\mathbf{X}$ 为列向量即退回向量二次型）；第三式在协方差/信息矩阵的最大似然里用到（\cref{ch:slam_est}）。本书已在点云配准（ICP 的 $\operatorname{tr}(\mathbf{R}\mathbf{W}^\top)$ 最大化）等处用到迹技巧。

\subsection{矩阵对矩阵求导}
\label{sec:matderiv-mat-mat}

\begin{note}
\textbf{矩阵对矩阵.}\;最一般的“矩阵对矩阵求导”需借助\emph{向量化算子} $\operatorname{vec}(\cdot)$（按列堆叠成长向量）与\emph{克罗内克积} $\otimes$，把它化归为“长向量对长向量”的雅可比，再套 \cref{def:matderiv-jacobian}。常用恒等式如 $\operatorname{vec}(\mathbf{A}\mathbf{X}\mathbf{B})=(\mathbf{B}^\top\otimes\mathbf{A})\operatorname{vec}(\mathbf{X})$\cite{magnus2019matrix}。这套机器在本书主线里很少直接出现——绝大多数“对矩阵的求导”都可经迹技巧 \cref{eq:matderiv-trace-id} 或李群扰动（\cref{ch:lie}，把对 $\mathbf{R}/\mathbf{T}$ 的求导转为对切空间向量 $\boldsymbol{\varphi}/\boldsymbol{\xi}$ 的求导）回避。故本节仅作索引，不展开；需要时可查 Magnus--Neudecker\cite{magnus2019matrix} 与 Matrix Cookbook\cite{petersen2012cookbook}。
\end{note}

\begin{insight}[与李群求导的关系]
为什么本书几乎不需要“矩阵对矩阵求导”？因为约束流形上的变量（$\mathbf{R}\in\SO(3)$、$\mathbf{T}\in\SE(3)$）\emph{不能}自由地逐元求导——逐元扰动会破坏正交性约束。\cref{ch:lie} 的做法是把对矩阵的求导\emph{转成}对切空间（李代数）向量的求导：$\partial(\mathbf{R}\mathbf{p})/\partial\boldsymbol{\varphi}=-(\mathbf{R}\mathbf{p})^\wedge$ 仍是一个标准的“向量对向量”雅可比（\cref{def:matderiv-jacobian}），只是分母 $\boldsymbol{\varphi}$ 住在 $\mathbb{R}^3$ 的切空间里。于是本附录的分子布局约定\emph{无缝}覆盖李群求导——这也是 \cref{ch:lie} 那条排布声明能直接套用的原因。
\end{insight}

% ════════════════════════════════════════════════════════════════
\section{矩阵代数补遗：vec/克罗内克积与分块恒等式（选读）}
\label{sec:matderiv-algebra}

上节把“矩阵对矩阵求导”归结到 $\operatorname{vec}$ 与克罗内克积 $\otimes$，但只点了一个恒等式作索引。这一组算子在本书少数几处\emph{绕不开}：协方差矩阵的最大似然要对 $\boldsymbol{\Sigma}$（对称矩阵）求导（\cref{ch:slam_est}）、连续时间先验要做矩阵指数（下节）、SLAM 信息阵的分块求逆与边缘化要用舒尔补（\cref{ch:nlopt}）。本节把这些\emph{代数}工具一次补齐并各给证明骨架，使全书在用到时不必外查。除非特别说明，下列矩阵维度都按使乘法相容来取，$\mathbf{1}$ 记单位阵。

\subsection{vec 算子与克罗内克积}
\label{sec:matderiv-vec-kron}

\begin{definition}[vec 算子与克罗内克积]\label{def:matderiv-vec-kron}
设 $\mathbf{A}=[\mathbf{a}_1\ \cdots\ \mathbf{a}_n]\in\mathbb{R}^{m\times n}$，其列为 $\mathbf{a}_j\in\mathbb{R}^m$。\emph{向量化算子}把各列\emph{自上而下}堆成一根长列：
\begin{equation}
  \operatorname{vec}(\mathbf{A})=\begin{bmatrix}\mathbf{a}_1\\\vdots\\\mathbf{a}_n\end{bmatrix}\in\mathbb{R}^{mn}.
  \label{eq:matderiv-vec}
\end{equation}
对 $\mathbf{A}\in\mathbb{R}^{m\times n}$、$\mathbf{B}\in\mathbb{R}^{p\times q}$，\emph{克罗内克积}（Kronecker product）是 $mp\times nq$ 的分块矩阵，把 $\mathbf{A}$ 的每个元 $A_{ij}$ 替换成块 $A_{ij}\mathbf{B}$：
\begin{equation}
  \mathbf{A}\otimes\mathbf{B}=\begin{bmatrix}
    A_{11}\mathbf{B}&\cdots&A_{1n}\mathbf{B}\\
    \vdots&\ddots&\vdots\\
    A_{m1}\mathbf{B}&\cdots&A_{mn}\mathbf{B}
  \end{bmatrix}\in\mathbb{R}^{mp\times nq}.
  \label{eq:matderiv-kron}
\end{equation}
\end{definition}

\begin{proposition}[vec--克罗内克核心恒等式]\label{prop:matderiv-vec-kron}
对维度相容的矩阵，有\cite{magnus2019matrix,barfoot2024state}
\begin{align}
  \operatorname{vec}(\mathbf{A}\mathbf{X}\mathbf{B})&=(\mathbf{B}^\top\otimes\mathbf{A})\operatorname{vec}(\mathbf{X}),
  \label{eq:matderiv-vec-AXB}\\
  \operatorname{vec}(\mathbf{a}\mathbf{b}^\top)&=\mathbf{b}\otimes\mathbf{a},\qquad
  \operatorname{vec}(\mathbf{A})^\top\operatorname{vec}(\mathbf{B})=\operatorname{tr}(\mathbf{A}^\top\mathbf{B}),
  \label{eq:matderiv-vec-tr}
\end{align}
且克罗内克积满足\emph{混合积律}与随之而来的转置、求逆、迹、行列式规则：
\begin{equation}
  (\mathbf{A}\otimes\mathbf{B})(\mathbf{C}\otimes\mathbf{D})=(\mathbf{A}\mathbf{C})\otimes(\mathbf{B}\mathbf{D}),
  \label{eq:matderiv-kron-mixed}
\end{equation}
\begin{equation}
  (\mathbf{A}\otimes\mathbf{B})^\top=\mathbf{A}^\top\otimes\mathbf{B}^\top,\quad
  (\mathbf{A}\otimes\mathbf{B})^{-1}=\mathbf{A}^{-1}\otimes\mathbf{B}^{-1},\quad
  \operatorname{tr}(\mathbf{A}\otimes\mathbf{B})=\operatorname{tr}(\mathbf{A})\operatorname{tr}(\mathbf{B}),
  \label{eq:matderiv-kron-props}
\end{equation}
对 $\mathbf{A}\in\mathbb{R}^{n\times n}$、$\mathbf{B}\in\mathbb{R}^{m\times m}$ 还有 $\det(\mathbf{A}\otimes\mathbf{B})=(\det\mathbf{A})^m(\det\mathbf{B})^n$。
\end{proposition}
\begin{proof}
核心式 \cref{eq:matderiv-vec-AXB}：记 $\mathbf{B}$ 的第 $j$ 列为 $\mathbf{b}_j$，则 $\mathbf{A}\mathbf{X}\mathbf{B}$ 的第 $j$ 列为 $\mathbf{A}\mathbf{X}\mathbf{b}_j=\mathbf{A}\mathbf{X}\sum_i B_{ij}\mathbf{1}_i=\sum_i B_{ij}\mathbf{A}\mathbf{x}_i$（$\mathbf{x}_i$ 为 $\mathbf{X}$ 的列），即 $(\mathbf{b}_j^\top\otimes\mathbf{A})\operatorname{vec}(\mathbf{X})$。把各 $j$ 列竖排即得 $(\mathbf{B}^\top\otimes\mathbf{A})\operatorname{vec}(\mathbf{X})$。\cref{eq:matderiv-vec-tr} 第二式按分量展开两边都等于 $\sum_{ij}A_{ij}B_{ij}$。混合积律 \cref{eq:matderiv-kron-mixed} 按 \cref{eq:matderiv-kron} 的分块逐块相乘 $(A_{ik}\mathbf{B})(C_{kj}\mathbf{D})$ 并对 $k$ 求和，正是 $((\mathbf{A}\mathbf{C})_{ij})(\mathbf{B}\mathbf{D})$；\cref{eq:matderiv-kron-props} 的求逆、迹由混合积律与定义直接得出。\hfill$\square$
\end{proof}

\begin{note}
\textbf{回填上节的索引.}\;\cref{eq:matderiv-vec-AXB} 正是 \cref{sec:matderiv-mat-mat} 那条“仅作索引”的 $\operatorname{vec}(\mathbf{A}\mathbf{X}\mathbf{B})=(\mathbf{B}^\top\otimes\mathbf{A})\operatorname{vec}(\mathbf{X})$。它把“矩阵对矩阵求导”机械化：要算 $\partial(\mathbf{A}\mathbf{X}\mathbf{B})/\partial\mathbf{X}$（按 $\operatorname{vec}$ 拉直后的雅可比），结果立刻读出为 $\mathbf{B}^\top\otimes\mathbf{A}$。本书主线仍优先用迹技巧（\cref{eq:matderiv-trace-id}）或李群扰动回避这套机器，但协方差求导、卡尔曼增益的灵敏度分析等处会直接调用它。
\end{note}

\subsection{对称矩阵的无冗余参数化与求导}
\label{sec:matderiv-vech}

协方差/信息矩阵是\emph{对称}的：$\boldsymbol{\Sigma}\in\mathbb{R}^{n\times n}$ 只有 $n(n+1)/2$ 个自由元，但 $\operatorname{vec}(\boldsymbol{\Sigma})$ 有 $n^2$ 个分量，把对角外的元重复计了两次。直接对 $\operatorname{vec}(\boldsymbol{\Sigma})$ 求导会得到“两份”导数；正确处理须显式扣除这重冗余。

\begin{definition}[半向量化与复制矩阵]\label{def:matderiv-vech}
\emph{半向量化算子} $\operatorname{vech}(\cdot)$ 只堆叠对称矩阵\emph{主对角及其下方}的元，得长度 $n(n+1)/2$ 的列。\emph{复制矩阵}（duplication matrix）$\mathbf{D}\in\mathbb{R}^{n^2\times n(n+1)/2}$ 把无冗余表示还原成全 $\operatorname{vec}$\cite{magnus2019matrix,barfoot2024state}：
\begin{equation}
  \operatorname{vec}(\boldsymbol{\Sigma})=\mathbf{D}\,\operatorname{vech}(\boldsymbol{\Sigma}),\qquad
  \operatorname{vech}(\boldsymbol{\Sigma})=\mathbf{D}^{+}\operatorname{vec}(\boldsymbol{\Sigma}),\quad \mathbf{D}^{+}=(\mathbf{D}^\top\mathbf{D})^{-1}\mathbf{D}^\top,
  \label{eq:matderiv-dup}
\end{equation}
其中 $\mathbf{D}^{+}$ 为 $\mathbf{D}$ 的 Moore--Penrose 伪逆，满足 $\mathbf{D}^{+}\mathbf{D}=\mathbf{1}$。
\end{definition}

\begin{proposition}[对称矩阵求导：对称化导数]\label{prop:matderiv-sym}
设 $f$ 是对称矩阵 $\mathbf{X}$ 的标量函数。把 $\mathbf{X}$ 暂当\emph{无约束}矩阵算出形式导数 $\partial f/\partial\mathbf{X}$（\cref{def:matderiv-scalar-mat}），则考虑对称约束后的“真导数”（记 $\bar{\partial} f/\bar{\partial}\mathbf{X}$，横杠提醒已计入对称性）是其\emph{对称化}：
\begin{equation}
  \frac{\bar{\partial} f}{\bar{\partial}\mathbf{X}}
  =\operatorname{sym}\!\Big(\frac{\partial f}{\partial\mathbf{X}}\Big)
  :=\frac{\partial f}{\partial\mathbf{X}}+\Big(\frac{\partial f}{\partial\mathbf{X}}\Big)^{\!\top}-\frac{\partial f}{\partial\mathbf{X}}\circ\mathbf{1},
  \label{eq:matderiv-sym}
\end{equation}
$\circ$ 为逐元（Hadamard）积；末项扣去被重复相加的主对角。
\end{proposition}
\begin{proof}
对对称 $\mathbf{X}$ 有 $\mathrm{d}\operatorname{vec}(\mathbf{X})=\mathbf{D}\,\mathrm{d}\operatorname{vech}(\mathbf{X})$。代入微分识别式（\cref{eq:matderiv-trace-id} 的 $\operatorname{vec}$ 写法）$\mathrm{d}f=\operatorname{vec}(\partial f/\partial\mathbf{X})^\top\mathrm{d}\operatorname{vec}(\mathbf{X})$，得 $\partial f/\partial\operatorname{vech}(\mathbf{X})^\top=\mathbf{D}^\top\operatorname{vec}(\partial f/\partial\mathbf{X})$；$\mathbf{D}^\top$ 的作用恰是把对角外“上、下两处”同一参数的偏导\emph{相加并归到唯一一份}。再用 $\mathbf{D}\mathbf{D}^\top\operatorname{vec}(\cdot)=\operatorname{vec}(\operatorname{sym}(\cdot))$ 还原回矩阵形，即 \cref{eq:matderiv-sym}。\hfill$\square$
\end{proof}

\begin{pitfall}[对协方差求导漏掉对称化]
最大似然估计协方差时，目标含 $\ln\det\boldsymbol{\Sigma}$ 与 $\operatorname{tr}(\boldsymbol{\Sigma}^{-1}\mathbf{S})$。把 $\boldsymbol{\Sigma}$ 当无约束矩阵求导会得到非对称的中间结果，\emph{若直接令其为零}会解出非对称的 $\boldsymbol{\Sigma}$。正确做法是按 \cref{prop:matderiv-sym} 取对称化导数再置零（或等价地，从一开始就只对 $\operatorname{vech}(\boldsymbol{\Sigma})$ 这 $n(n+1)/2$ 个独立参数求导）。所幸 $\ln\det$ 与迹这两项的无约束导数 $\boldsymbol{\Sigma}^{-\top}$、$-\boldsymbol{\Sigma}^{-\top}\mathbf{S}\boldsymbol{\Sigma}^{-\top}$ 本就对称，对称化只改对角缩放、不改“令导数为零”的解 $\hat{\boldsymbol{\Sigma}}=\mathbf{S}$——这又是“对称帮了忙”，不是普遍规律。
\end{pitfall}

\subsection{分块矩阵求逆与舒尔补}
\label{sec:matderiv-block-inv}

\begin{proposition}[分块求逆与舒尔补]\label{prop:matderiv-block-inv}
设分块矩阵 $\mathbf{M}=\big[\begin{smallmatrix}\mathbf{A}&\mathbf{B}\\\mathbf{C}&\mathbf{D}\end{smallmatrix}\big]$，$\mathbf{A},\mathbf{D}$ 方阵且可逆。记 $\mathbf{D}$ 的\emph{舒尔补} $\mathbf{M}/\mathbf{D}:=\mathbf{A}-\mathbf{B}\mathbf{D}^{-1}\mathbf{C}$。当 $\mathbf{M}/\mathbf{D}$ 可逆时\cite{petersen2012cookbook,barfoot2024state}
\begin{equation}
  \mathbf{M}^{-1}=
  \begin{bmatrix}
    (\mathbf{M}/\mathbf{D})^{-1} & -(\mathbf{M}/\mathbf{D})^{-1}\mathbf{B}\mathbf{D}^{-1}\\
    -\mathbf{D}^{-1}\mathbf{C}(\mathbf{M}/\mathbf{D})^{-1} & \mathbf{D}^{-1}+\mathbf{D}^{-1}\mathbf{C}(\mathbf{M}/\mathbf{D})^{-1}\mathbf{B}\mathbf{D}^{-1}
  \end{bmatrix},
  \label{eq:matderiv-block-inv}
\end{equation}
且行列式按舒尔补\emph{相乘分解}：$\det\mathbf{M}=\det\mathbf{D}\,\det(\mathbf{M}/\mathbf{D})=\det\mathbf{A}\,\det(\mathbf{M}/\mathbf{A})$，其中 $\mathbf{M}/\mathbf{A}=\mathbf{D}-\mathbf{C}\mathbf{A}^{-1}\mathbf{B}$。
\end{proposition}
\begin{proof}
对 $\mathbf{M}$ 作分块 LDU 分解
\begin{equation*}
  \mathbf{M}=\begin{bmatrix}\mathbf{1}&\mathbf{B}\mathbf{D}^{-1}\\\mathbf{0}&\mathbf{1}\end{bmatrix}
  \begin{bmatrix}\mathbf{M}/\mathbf{D}&\mathbf{0}\\\mathbf{0}&\mathbf{D}\end{bmatrix}
  \begin{bmatrix}\mathbf{1}&\mathbf{0}\\\mathbf{D}^{-1}\mathbf{C}&\mathbf{1}\end{bmatrix},
\end{equation*}
三因子均可逆，逐个求逆（外两块单位三角阵的逆只需翻号）并相乘即 \cref{eq:matderiv-block-inv}；对三因子取行列式，单位三角阵行列式为 $1$，中间块对角，得 $\det\mathbf{M}=\det(\mathbf{M}/\mathbf{D})\det\mathbf{D}$。换成 UDL 分解（先消左下）得 $\mathbf{M}/\mathbf{A}$ 版本。\hfill$\square$
\end{proof}

\begin{insight}[舒尔补：边缘化与条件化的同一根源]
\cref{prop:matderiv-block-inv} 是本书反复出现的代数引擎。\textbf{(1) 边缘化}：联合高斯 $\big[\begin{smallmatrix}\mathbf{x}\\\mathbf{y}\end{smallmatrix}\big]$ 的信息阵作 $\mathbf{y}$ 块的舒尔补，得 $\mathbf{x}$ 的\emph{边缘}信息阵——这正是 \cref{ch:slam_est} 把路标边缘化、\cref{ch:nlopt} 求解箭头型 BA/SLAM 法方程的核心（先消路标块再解位姿）。\textbf{(2) 求逆引理}：\cref{eq:matderiv-block-inv} 左上块 $(\mathbf{A}-\mathbf{B}\mathbf{D}^{-1}\mathbf{C})^{-1}$ 与右下块对照，立得 \cref{ch:slam_est} 的 Sherman--Morrison--Woodbury 引理。\textbf{(3) 行列式分解}：用于高斯归一化常数与互信息（\cref{ch:slam_est} 的 D-最优主动感知）。这与 \cref{ch:slam_est} 用 LDU/UDL 比较块得 SMW 是\emph{同一手法}的两面。
\end{insight}

% ════════════════════════════════════════════════════════════════
\section{Van Loan 分块矩阵指数：连续 LTI 系统的离散化（选读）}
\label{sec:matderiv-vanloan}

本节给出一个在状态估计里极其实用、却常被当“黑箱”的结果：用\emph{一次}矩阵指数同时算出连续时间线性时不变（LTI）随机微分方程（SDE）离散化所需的\emph{状态转移矩阵} $\boldsymbol{\Phi}$ \emph{与}\emph{离散过程噪声协方差} $\mathbf{Q}_k$。这是 Van Loan（1978）的经典技巧\cite{vanloan1978computing}。它是 \cref{ch:eskf} 的连续$\to$离散噪声换算、\cref{ch:slam_est} 连续时间轨迹估计（STEAM，\cref{sec:est-steam}）里 WNOA 先验的 $\boldsymbol{\Phi},\mathbf{Q}$ 的统一来源——那里因系统矩阵\emph{幂零}而能手算，本节给的是\emph{任意}LTI 系统都适用的机械算法。

\subsection{问题：连续 LTI SDE 的离散化}
\label{sec:matderiv-vanloan-setup}

\begin{definition}[连续 LTI SDE 与其离散等价]\label{def:matderiv-lti}
考虑连续时间 LTI SDE
\begin{equation}
  \dot{\mathbf{x}}(t)=\mathbf{A}\,\mathbf{x}(t)+\mathbf{L}\,\mathbf{w}(t),
  \qquad \mathbf{w}(t)\sim\mathcal{GP}\big(\mathbf{0},\,\mathbf{Q}_C\,\delta(t-t')\big),
  \label{eq:matderiv-lti-sde}
\end{equation}
其中 $\mathbf{A}\in\mathbb{R}^{n\times n}$、$\mathbf{L}\in\mathbb{R}^{n\times l}$ 为常矩阵，$\mathbf{Q}_C\in\mathbb{R}^{l\times l}$ 为白噪声的\emph{功率谱密度}（PSD）。在固定步长 $\Delta t$ 下采样，其精确离散等价为线性高斯过程
\begin{equation}
  \mathbf{x}_k=\boldsymbol{\Phi}\,\mathbf{x}_{k-1}+\mathbf{w}_{k},
  \qquad \mathbf{w}_{k}\sim\mathcal{N}(\mathbf{0},\,\mathbf{Q}_k),
  \label{eq:matderiv-lti-disc}
\end{equation}
其中状态转移矩阵与离散过程噪声协方差为
\begin{equation}
  \boldsymbol{\Phi}=\exp(\mathbf{A}\,\Delta t)\in\mathbb{R}^{n\times n},
  \qquad
  \mathbf{Q}_k=\int_{0}^{\Delta t}\!\exp(\mathbf{A}s)\,\mathbf{L}\mathbf{Q}_C\mathbf{L}^\top\,\exp(\mathbf{A}s)^\top\,\mathrm{d}s\in\mathbb{R}^{n\times n}.
  \label{eq:matderiv-PhiQ-def}
\end{equation}
\end{definition}

\begin{note}
\textbf{两式从何而来.}\;转移矩阵 $\boldsymbol{\Phi}(t,s)$ 解齐次方程 $\dot{\boldsymbol{\Phi}}=\mathbf{A}\boldsymbol{\Phi}$、$\boldsymbol{\Phi}(s,s)=\mathbf{1}$；LTI 下只依赖时差，$\boldsymbol{\Phi}=\exp(\mathbf{A}\,\Delta t)$。把 \cref{eq:matderiv-lti-sde} 的解 $\mathbf{x}(t_k)=\boldsymbol{\Phi}\,\mathbf{x}(t_{k-1})+\int_0^{\Delta t}\exp(\mathbf{A}s)\mathbf{L}\,\mathbf{w}(t_k-s)\,\mathrm{d}s$ 中噪声项的协方差用白噪声的 $\delta$ 相关性 $\mathbb{E}[\mathbf{w}(\tau)\mathbf{w}(\tau')^\top]=\mathbf{Q}_C\delta(\tau-\tau')$ 算出，即得 \cref{eq:matderiv-PhiQ-def} 的 $\mathbf{Q}_k$ 积分。难点在这个\emph{矩阵积分}：被积函数含 $\exp(\mathbf{A}s)$ 与其转置的夹心，一般\emph{没有}简单闭式。
\end{note}

\subsection{Van Loan 技巧：一次矩阵指数算齐 \texorpdfstring{$\boldsymbol{\Phi}$}{Phi} 与 \texorpdfstring{$\mathbf{Q}_k$}{Qk}}
\label{sec:matderiv-vanloan-trick}

\begin{theorem}[Van Loan 1978]\label{thm:matderiv-vanloan}
构造 $2n\times2n$ 的\emph{增广矩阵}
\begin{equation}
  \boldsymbol{\Xi}=\begin{bmatrix}-\mathbf{A}&\mathbf{L}\mathbf{Q}_C\mathbf{L}^\top\\[2pt]\mathbf{0}&\mathbf{A}^\top\end{bmatrix}\in\mathbb{R}^{2n\times2n},
  \label{eq:matderiv-vanloan-aug}
\end{equation}
并对它取一次矩阵指数，把结果按 $n\times n$ 分块记为
\begin{equation}
  \exp(\boldsymbol{\Xi}\,\Delta t)
  =\begin{bmatrix}\boldsymbol{\Upsilon}_{11}&\boldsymbol{\Upsilon}_{12}\\\mathbf{0}&\boldsymbol{\Upsilon}_{22}\end{bmatrix}.
  \label{eq:matderiv-vanloan-exp}
\end{equation}
则 \cref{def:matderiv-lti} 所需的两块可直接\emph{读出}：
\begin{equation}
  \boxed{\ \boldsymbol{\Phi}=\boldsymbol{\Upsilon}_{22}^\top,\qquad
  \mathbf{Q}_k=\boldsymbol{\Upsilon}_{22}^\top\,\boldsymbol{\Upsilon}_{12}\ }
  \label{eq:matderiv-vanloan-readout}
\end{equation}
（等价地 $\boldsymbol{\Phi}=\exp(\mathbf{A}\Delta t)$ 由右下块转置给出，$\mathbf{Q}_k=\boldsymbol{\Phi}\,\boldsymbol{\Upsilon}_{12}$）。
\end{theorem}
\begin{proof}
因 $\boldsymbol{\Xi}$ 是分块上三角，其指数也分块上三角，且对角块各自为子块的指数：
\begin{equation*}
  \boldsymbol{\Upsilon}_{11}=\exp(-\mathbf{A}\,\Delta t),\qquad
  \boldsymbol{\Upsilon}_{22}=\exp(\mathbf{A}^\top\Delta t)=\exp(\mathbf{A}\,\Delta t)^\top=\boldsymbol{\Phi}^\top,
\end{equation*}
故 $\boldsymbol{\Upsilon}_{22}^\top=\boldsymbol{\Phi}$，这给出 \cref{eq:matderiv-vanloan-readout} 左式。对右上块 $\boldsymbol{\Upsilon}_{12}$，令 $\mathbf{U}(t)=\exp(\boldsymbol{\Xi}t)$，由 $\dot{\mathbf{U}}=\boldsymbol{\Xi}\mathbf{U}$ 取 $(1,2)$ 块得线性 ODE
\begin{equation*}
  \dot{\boldsymbol{\Upsilon}}_{12}(t)=-\mathbf{A}\,\boldsymbol{\Upsilon}_{12}(t)+\mathbf{L}\mathbf{Q}_C\mathbf{L}^\top\,\boldsymbol{\Upsilon}_{22}(t),\qquad \boldsymbol{\Upsilon}_{12}(0)=\mathbf{0}.
\end{equation*}
代入已知的 $\boldsymbol{\Upsilon}_{22}(t)=\exp(\mathbf{A}^\top t)$，用积分因子 $\exp(\mathbf{A}t)$ 解这条一阶线性矩阵 ODE：
\begin{equation*}
  \boldsymbol{\Upsilon}_{12}(\Delta t)=\int_0^{\Delta t}\exp(-\mathbf{A}(\Delta t-s))\,\mathbf{L}\mathbf{Q}_C\mathbf{L}^\top\,\exp(\mathbf{A}^\top s)\,\mathrm{d}s.
\end{equation*}
左乘 $\boldsymbol{\Phi}=\exp(\mathbf{A}\,\Delta t)$ 并把 $\exp(\mathbf{A}\,\Delta t)\exp(-\mathbf{A}(\Delta t-s))=\exp(\mathbf{A}s)$ 合并：
\begin{equation*}
  \boldsymbol{\Phi}\,\boldsymbol{\Upsilon}_{12}(\Delta t)
  =\int_0^{\Delta t}\exp(\mathbf{A}s)\,\mathbf{L}\mathbf{Q}_C\mathbf{L}^\top\,\exp(\mathbf{A}^\top s)\,\mathrm{d}s,
\end{equation*}
右端恰是 \cref{eq:matderiv-PhiQ-def} 定义的 $\mathbf{Q}_k$（注意 $\exp(\mathbf{A}^\top s)=\exp(\mathbf{A}s)^\top$）。故 $\mathbf{Q}_k=\boldsymbol{\Phi}\,\boldsymbol{\Upsilon}_{12}=\boldsymbol{\Upsilon}_{22}^\top\,\boldsymbol{\Upsilon}_{12}$，即 \cref{eq:matderiv-vanloan-readout} 右式。\hfill$\square$
\end{proof}

\begin{insight}[为什么这是“一次算齐”]
\cref{thm:matderiv-vanloan} 的妙处在于：把那个棘手的\emph{矩阵积分} \cref{eq:matderiv-PhiQ-def} 转化为一个纯\emph{矩阵指数}，而后者有成熟稳定的数值算法（如 scaling-and-squaring）。增广矩阵右下块的 $\mathbf{A}^\top$ 负责生成 $\boldsymbol{\Phi}^\top$、左上块的 $-\mathbf{A}$ 与右上块的 $\mathbf{L}\mathbf{Q}_C\mathbf{L}^\top$ 协同把积分“积”进 $\boldsymbol{\Upsilon}_{12}$——三块各司其职，一次指数把 $\boldsymbol{\Phi}$ 与 $\mathbf{Q}_k$ 同时吐出。工程上这是 EKF/IMU 预积分（\cref{ch:eskf}、IMU 章）和连续时间估计实现里离散化的\emph{标准引擎}。
\end{insight}

\begin{derivation}[校核：用 Van Loan 复算 STEAM 的 WNOA 先验]
\cref{sec:est-steam} 的 WNOA（白噪声加速度）先验取 $\mathbf{A}=\big[\begin{smallmatrix}\mathbf{0}&\mathbf{1}\\\mathbf{0}&\mathbf{0}\end{smallmatrix}\big]$、$\mathbf{L}=\big[\begin{smallmatrix}\mathbf{0}\\\mathbf{1}\end{smallmatrix}\big]$（每块 $6\times6$，PSD 记 $\mathbf{Q}_C$），那里因 $\mathbf{A}^2=\mathbf{0}$ \emph{幂零}而手算出 \cref{eq:steam-Phi-Q}。这里用 \cref{thm:matderiv-vanloan} 机械复现，验证两法一致。先算各子块指数：$\boldsymbol{\Phi}=\exp(\mathbf{A}\Delta t)=\mathbf{1}+\mathbf{A}\Delta t=\big[\begin{smallmatrix}\mathbf{1}&\Delta t\,\mathbf{1}\\\mathbf{0}&\mathbf{1}\end{smallmatrix}\big]$（与 \cref{eq:steam-Phi-Q} 左式一致）。再代 \cref{eq:matderiv-PhiQ-def} 的积分（与解 $\boldsymbol{\Upsilon}_{12}$ 等价）：$\exp(\mathbf{A}s)\mathbf{L}=[s\,\mathbf{1};\,\mathbf{1}]$，故
\[
  \mathbf{Q}_k=\int_0^{\Delta t}\begin{bmatrix}s\,\mathbf{1}\\\mathbf{1}\end{bmatrix}\mathbf{Q}_C\begin{bmatrix}s\,\mathbf{1}&\mathbf{1}\end{bmatrix}\mathrm{d}s
  =\int_0^{\Delta t}\begin{bmatrix}s^2\mathbf{Q}_C&s\,\mathbf{Q}_C\\s\,\mathbf{Q}_C&\mathbf{Q}_C\end{bmatrix}\mathrm{d}s
  =\begin{bmatrix}\tfrac13\Delta t^3\mathbf{Q}_C&\tfrac12\Delta t^2\mathbf{Q}_C\\[2pt]\tfrac12\Delta t^2\mathbf{Q}_C&\Delta t\,\mathbf{Q}_C\end{bmatrix},
\]
正是 \cref{eq:steam-Phi-Q} 右式。\textbf{要点}：WNOA 之所以能\emph{手算}，是 $\mathbf{A}$ 幂零这一特殊性；Van Loan 不需要这个前提——换成 WNOJ（白噪声加加速度，$\mathbf{A}$ 仍幂零但更大）、Singer 模型（$\mathbf{A}$ 含\emph{非零}衰减项、\emph{不再}幂零）等更一般先验，手算积分会很繁，而 \cref{thm:matderiv-vanloan} 一律适用。这正是把 \cref{sec:est-steam} 那条“换 SDE 只换 $\mathbf{A},\mathbf{L},\boldsymbol{\Phi},\mathbf{Q}$”落到\emph{可计算}的统一离散化算法。\hfill$\square$
\end{derivation}

\begin{note}
\textbf{含控制输入的推广.}\;若 SDE 带确定性输入 $\dot{\mathbf{x}}=\mathbf{A}\mathbf{x}+\mathbf{B}\mathbf{u}+\mathbf{L}\mathbf{w}$（$\mathbf{u}$ 在步内为常），离散输入矩阵 $\mathbf{B}_k=\int_0^{\Delta t}\exp(\mathbf{A}s)\,\mathrm{d}s\,\mathbf{B}$ 可用同一思路：对 $\big[\begin{smallmatrix}\mathbf{A}&\mathbf{B}\\\mathbf{0}&\mathbf{0}\end{smallmatrix}\big]\Delta t$ 取指数，右上块即 $\int_0^{\Delta t}\exp(\mathbf{A}s)\mathrm{d}s\,\mathbf{B}=\mathbf{B}_k$。Van Loan 的原文给出同时含 $\boldsymbol{\Phi}$、$\mathbf{B}_k$、$\mathbf{Q}_k$ 乃至积分二次型的更大增广矩阵\cite{vanloan1978computing}；本书只需 $\boldsymbol{\Phi},\mathbf{Q}_k$ 这一最常用情形。
\end{note}

% ════════════════════════════════════════════════════════════════
\section{Fisher 信息阵在重参数化下的变换（选读）}
\label{sec:matderiv-fim}

\cref{ch:slam_est} 的 \cref{thm:crlb} 给出 Fisher 信息阵（FIM）$\boldsymbol{\mathcal{I}}_{\boldsymbol{\theta}}$ 与克拉默--拉奥下界（CRLB）$\mathrm{cov}(\hat{\boldsymbol{\theta}})\succeq\boldsymbol{\mathcal{I}}_{\boldsymbol{\theta}}^{-1}$。一个常见而要紧的问题：同一统计模型，换一组参数（如把协方差换成信息形式、把欧拉角换成四元数、把 $\SE(3)$ 位姿换成李代数坐标）后，FIM 与 CRLB 如何随之变换？答案是一个干净的\emph{雅可比夹心}，本质就是本附录的链式法则（\cref{thm:matderiv-chain}）作用在“对数似然的曲率”上。

\begin{theorem}[FIM 的重参数化变换]\label{thm:matderiv-fim-reparam}
设统计模型在参数 $\boldsymbol{\theta}\in\mathbb{R}^d$ 下的 FIM 为 $\boldsymbol{\mathcal{I}}_{\boldsymbol{\theta}}$。换一组参数 $\boldsymbol{\phi}\in\mathbb{R}^{d}$，二者由可微\emph{双射} $\boldsymbol{\theta}=\mathbf{g}(\boldsymbol{\phi})$ 联系，记雅可比
\begin{equation}
  \mathbf{J}=\frac{\partial\boldsymbol{\theta}}{\partial\boldsymbol{\phi}^\top}=\frac{\partial\mathbf{g}}{\partial\boldsymbol{\phi}^\top}\in\mathbb{R}^{d\times d}\quad(\text{分子布局，\cref{def:matderiv-jacobian}}).
  \label{eq:matderiv-fim-J}
\end{equation}
则新参数下的 FIM 是旧 FIM 的雅可比夹心：
\begin{equation}
  \boxed{\ \boldsymbol{\mathcal{I}}_{\boldsymbol{\phi}}=\mathbf{J}^\top\,\boldsymbol{\mathcal{I}}_{\boldsymbol{\theta}}\,\mathbf{J}\ }.
  \label{eq:matderiv-fim-reparam}
\end{equation}
\end{theorem}
\begin{proof}
记对数似然 $\ell(\boldsymbol{\theta})=\ln p(\mathbf{x}\mid\boldsymbol{\theta})$，\emph{得分}（score）按本书梯度约定取列向量 $\mathbf{s}_{\boldsymbol{\theta}}=\partial\ell/\partial\boldsymbol{\theta}$。FIM 的一个等价定义是得分的协方差 $\boldsymbol{\mathcal{I}}_{\boldsymbol{\theta}}=\mathbb{E}[\mathbf{s}_{\boldsymbol{\theta}}\mathbf{s}_{\boldsymbol{\theta}}^\top]$（在正则条件下等于 \cref{eq:crlb} 的期望海森）。复合 $\ell(\mathbf{g}(\boldsymbol{\phi}))$，由链式法则（\cref{thm:matderiv-chain}）得分按列向量变换为
\begin{equation*}
  \mathbf{s}_{\boldsymbol{\phi}}=\frac{\partial\ell}{\partial\boldsymbol{\phi}}
  =\Big(\frac{\partial\boldsymbol{\theta}}{\partial\boldsymbol{\phi}^\top}\Big)^{\!\top}\frac{\partial\ell}{\partial\boldsymbol{\theta}}
  =\mathbf{J}^\top\mathbf{s}_{\boldsymbol{\theta}}.
\end{equation*}
（标量对列向量求导得列、且复合时外层雅可比转置左乘，见 \cref{prop:matderiv-transpose} 与 \cref{app:matderiv-layout-table}。）代入协方差定义：
\begin{equation*}
  \boldsymbol{\mathcal{I}}_{\boldsymbol{\phi}}=\mathbb{E}[\mathbf{s}_{\boldsymbol{\phi}}\mathbf{s}_{\boldsymbol{\phi}}^\top]
  =\mathbb{E}[\mathbf{J}^\top\mathbf{s}_{\boldsymbol{\theta}}\mathbf{s}_{\boldsymbol{\theta}}^\top\mathbf{J}]
  =\mathbf{J}^\top\mathbb{E}[\mathbf{s}_{\boldsymbol{\theta}}\mathbf{s}_{\boldsymbol{\theta}}^\top]\mathbf{J}
  =\mathbf{J}^\top\boldsymbol{\mathcal{I}}_{\boldsymbol{\theta}}\mathbf{J},
\end{equation*}
其中 $\mathbf{J}$ 是确定性的（不依赖 $\mathbf{x}$），可提出期望。\hfill$\square$
\end{proof}

\begin{corollary}[CRLB 的相应变换]\label{cor:matderiv-crlb-reparam}
当 $\mathbf{J}$ 可逆时，\cref{eq:matderiv-fim-reparam} 给逆 FIM（即 CRLB 下界）的变换
\begin{equation}
  \boldsymbol{\mathcal{I}}_{\boldsymbol{\phi}}^{-1}=\mathbf{J}^{-1}\,\boldsymbol{\mathcal{I}}_{\boldsymbol{\theta}}^{-1}\,\mathbf{J}^{-\top}.
  \label{eq:matderiv-crlb-reparam}
\end{equation}
这与无偏估计的协方差经一阶 delta 法 $\mathrm{cov}(\hat{\boldsymbol{\phi}})\approx\mathbf{J}^{-1}\mathrm{cov}(\hat{\boldsymbol{\theta}})\mathbf{J}^{-\top}$ 的传播\emph{形状完全一致}（\cref{ch:slam_est} 协方差传播 $\mathbf{F}\boldsymbol{\Sigma}\mathbf{F}^\top$ 的特例），故“达到 CRLB”这一性质在可逆重参数化下\emph{不变}：有效估计器换参数后仍有效。
\end{corollary}

\begin{insight}[为什么本书关心这件事]
三处直接用到 \cref{thm:matderiv-fim-reparam}。\textbf{(1) 协方差 vs 信息形式}：FIM 在“均值$+$协方差”参数化与“信息向量$+$信息阵”参数化下不同，按 \cref{eq:matderiv-fim-reparam} 互换（\cref{ch:slam_est} 的协方差$\leftrightarrow$信息对偶）。\textbf{(2) 流形上的 CRLB}：$\SO(3)/\SE(3)$ 上估计时，把 FIM 从过参数化（如 $9$ 个旋转元）拉回\emph{李代数}最小坐标（$3$ 或 $6$ 维），用的正是 \cref{eq:matderiv-fim-reparam}，$\mathbf{J}$ 取指数映射的雅可比——这让 CRLB 定义在切空间上、与 \cref{ch:lie} 的扰动求导自洽。\textbf{(3) 量纲与条件数}：换参数（如对数尺度）可改善 FIM 的条件数；\cref{eq:matderiv-fim-reparam} 量化了这一改变。注意若 $\mathbf{g}$ 是\emph{升维}嵌入（$\boldsymbol{\phi}$ 维数 $>$ $\boldsymbol{\theta}$），$\mathbf{J}$ 非方，\cref{eq:matderiv-fim-reparam} 仍成立但 $\boldsymbol{\mathcal{I}}_{\boldsymbol{\phi}}$ 奇异（过参数化的规范自由度，对应 \cref{ch:slam_est} 的 gauge 零空间），此时 CRLB 须用伪逆。
\end{insight}

% ════════════════════════════════════════════════════════════════
\section{小结与速查}
\label{sec:matderiv-summary}

\paragraph{一图看清两种布局。}
\cref{app:matderiv-layout-fig} 把分子布局与分母布局的形状对照画出来：同一组偏导，分子布局摆成 $n\times m$（行随分子、列随分母），分母布局摆成其转置 $m\times n$。本书全程走左边。

\begin{figure}[H]\centering
\begin{tikzpicture}[font=\small,>=stealth]
  % numerator layout block (book convention)
  \begin{scope}[shift={(0,0)}]
    \draw[thick,fill=structurecolor!8] (0,0) rectangle (3.0,2.0);
    \node at (1.5,2.35) {分子布局（本书）$\dfrac{\partial\mathbf{F}}{\partial\mathbf{x}^\top}\in\mathbb{R}^{n\times m}$};
    % row arrow (numerator F, n rows)
    \draw[->,thick] (-0.35,2.0) -- (-0.35,0.0) node[midway,left,align=center]{行 $\,n$\\(随分子 $\mathbf{F}$)};
    % col arrow (denominator x, m cols)
    \draw[->,thick] (0.0,-0.35) -- (3.0,-0.35) node[midway,below,align=center]{列 $\,m$（随分母 $\mathbf{x}$）};
    \node at (1.5,1.0) {$[\,\partial F_k/\partial x_l\,]$};
  \end{scope}
  % transpose relation
  \node at (4.6,1.0) {$\xrightarrow{\ (\cdot)^\top\ }$};
  % denominator layout block
  \begin{scope}[shift={(6.2,0.25)}]
    \draw[thick,fill=second!8] (0,0) rectangle (2.0,3.0);
    \node at (1.0,3.35) {分母布局 $\dfrac{\partial\mathbf{F}}{\partial\mathbf{x}}\in\mathbb{R}^{m\times n}$};
    \draw[->,thick] (-0.35,3.0) -- (-0.35,0.0) node[midway,left,align=center]{行 $\,m$};
    \draw[->,thick] (0.0,-0.35) -- (2.0,-0.35) node[midway,below,align=center]{列 $\,n$};
    \node at (1.0,1.5) {$[\,\partial F_k/\partial x_l\,]^\top$};
  \end{scope}
\end{tikzpicture}
\caption{分子布局（左，$n\times m$，本书统一采用）与分母布局（右，$m\times n$）互为转置。行随分子、列随分母是分子布局的记忆口诀，与 \cref{ch:lie} 的扰动雅可比排布、\cref{ch:nlopt} 的残差雅可比 $\partial\mathbf{e}/\partial\mathbf{x}$（行$=$残差维、列$=$状态维）一致。}
\label{app:matderiv-layout-fig}
\end{figure}

\paragraph{核心结论速查。}
\begin{itemize}
  \item \textbf{定义}：标量对向量 $=$ 梯度（\cref{def:matderiv-grad}，本书取列向量 $\nabla f$）；向量对向量 $=$ 雅可比 $n\times m$（\cref{def:matderiv-jacobian}）。
  \item \textbf{约定}：全书\emph{分子布局}（行随分子、列随分母，\cref{def:matderiv-layout}），与 \cref{ch:lie}、十四讲、Ceres/g2o/GTSAM 一致；简记可省分母转置号但不改形状（\cref{eq:matderiv-shorthand}）。
  \item \textbf{公式}：见 \cref{app:matderiv-formula-table}（线性型、内积、平方、二次型、链式、复合线性）。
  \item \textbf{贯穿例}：$\partial\|\mathbf{A}\mathbf{x}-\mathbf{b}\|^2/\partial\mathbf{x}=\mathbf{0}\Rightarrow\mathbf{A}^\top\mathbf{A}\mathbf{x}=\mathbf{A}^\top\mathbf{b}$（\cref{eq:matderiv-normal}），即最小二乘正规方程，全书估计的代数骨架。
  \item \textbf{延伸（选读）}：标量对矩阵用迹技巧 $\mathrm{d}f=\operatorname{tr}(\mathbf{G}^\top\mathrm{d}\mathbf{X})\Rightarrow\partial f/\partial\mathbf{X}=\mathbf{G}$（\cref{eq:matderiv-trace-id}）；矩阵对矩阵借 $\operatorname{vec}$ 与 $\otimes$，本书多用李群扰动回避。
  \item \textbf{代数补遗（选读）}：vec/克罗内克核心式 $\operatorname{vec}(\mathbf{A}\mathbf{X}\mathbf{B})=(\mathbf{B}^\top\otimes\mathbf{A})\operatorname{vec}(\mathbf{X})$（\cref{prop:matderiv-vec-kron}）；对称矩阵求导取对称化导数（\cref{prop:matderiv-sym}）；分块求逆/舒尔补（\cref{prop:matderiv-block-inv}）统一了边缘化、SMW、行列式分解。
  \item \textbf{离散化引擎（选读）}：Van Loan 增广矩阵指数一次算齐连续 LTI SDE 的 $\boldsymbol{\Phi},\mathbf{Q}_k$（\cref{thm:matderiv-vanloan}），是 \cref{ch:eskf}、IMU、STEAM（\cref{sec:est-steam}）离散化的公共算法。
  \item \textbf{统计变换（选读）}：FIM 在重参数化 $\boldsymbol{\theta}=\mathbf{g}(\boldsymbol{\phi})$ 下作雅可比夹心 $\boldsymbol{\mathcal{I}}_{\boldsymbol{\phi}}=\mathbf{J}^\top\boldsymbol{\mathcal{I}}_{\boldsymbol{\theta}}\mathbf{J}$（\cref{thm:matderiv-fim-reparam}），CRLB 随之 $\mathbf{J}^{-1}(\cdot)\mathbf{J}^{-\top}$，达界性不变。
\end{itemize}

\paragraph{后续关系。}
本附录是全书雅可比的统一地基，被以下各章直接引用：\cref{ch:rigid_body}（坐标变换与运动的求导）、\cref{ch:lie}（右扰动雅可比 $\partial(\mathbf{R}\mathbf{p})/\partial\boldsymbol{\varphi}$ 与分块排布的依据）、\cref{ch:slam_est}（线性高斯估计、协方差传播 $\mathbf{F}\boldsymbol{\Sigma}\mathbf{F}^\top$、舒尔补边缘化、Fisher 信息与 CRLB、最大似然对协方差求导）、\cref{ch:eskf}（连续$\to$离散噪声离散化用 \cref{thm:matderiv-vanloan}）、\cref{ch:nlopt}（残差线性化、高斯牛顿正规方程 \cref{eq:gn-normal}、白化与马氏二次型求导）。其中 \cref{sec:matderiv-vanloan} 的 Van Loan 技巧为 \cref{sec:est-steam} 连续时间轨迹估计的 WNOA 先验 $\boldsymbol{\Phi},\mathbf{Q}$ 提供了\emph{不依赖幂零性}的统一推导。读这些章遇到“对某量求雅可比”时，回到 \cref{app:matderiv-formula-table} 即可对号入座。

\begin{practice}
\begin{enumerate}
  \item （定义）写出 $\mathbf{F}(\mathbf{x})=[x_1^2+x_2,\ \sin x_1,\ x_1 x_3]^\top$（$\mathbf{x}\in\mathbb{R}^3$）的雅可比 $\partial\mathbf{F}/\partial\mathbf{x}^\top\in\mathbb{R}^{3\times3}$，并核对每个 $(k,l)$ 元的行列归属（\cref{def:matderiv-jacobian}）。
  \item （公式）用 \cref{prop:matderiv-xtAx} 求 $\partial(\mathbf{x}^\top\mathbf{A}\mathbf{x})/\partial\mathbf{x}$，分别取 $\mathbf{A}=\big[\begin{smallmatrix}1&2\\0&3\end{smallmatrix}\big]$（非对称）与其对称化 $\tfrac12(\mathbf{A}+\mathbf{A}^\top)$，验证两者给出\emph{相同}梯度——解释为什么二次型只“看到”$\mathbf{A}$ 的对称部分。
  \item （链式）残差 $\mathbf{e}(\mathbf{x})=\mathbf{z}-\mathbf{A}\,\mathbf{g}(\mathbf{x})$，$\mathbf{g}:\mathbb{R}^m\to\mathbb{R}^p$ 可微、$\mathbf{A}\in\mathbb{R}^{n\times p}$ 常量。用 \cref{thm:matderiv-chain} 与 \cref{cor:matderiv-Au} 写出 $\partial\mathbf{e}/\partial\mathbf{x}^\top$，并说明它如何进入 \cref{ch:nlopt} 的正规方程 \cref{eq:gn-normal}。
  \item （贯穿例）把 \cref{eq:matderiv-ls-obj} 的目标改为加权 $\|\mathbf{A}\mathbf{x}-\mathbf{b}\|_{\boldsymbol{\Sigma}^{-1}}^2=(\mathbf{A}\mathbf{x}-\mathbf{b})^\top\boldsymbol{\Sigma}^{-1}(\mathbf{A}\mathbf{x}-\mathbf{b})$（$\boldsymbol{\Sigma}^{-1}$ 对称正定），重做三步推导，得加权正规方程 $\mathbf{A}^\top\boldsymbol{\Sigma}^{-1}\mathbf{A}\mathbf{x}=\mathbf{A}^\top\boldsymbol{\Sigma}^{-1}\mathbf{b}$，并与 \cref{ch:nlopt} 的 \cref{eq:nlopt-gls} 对照。
  \item （选读）用迹技巧 \cref{eq:matderiv-trace-id} 推 \cref{eq:matderiv-trace-formulas} 的第一式 $\partial\operatorname{tr}(\mathbf{A}\mathbf{X})/\partial\mathbf{X}=\mathbf{A}^\top$（提示：$\mathrm{d}\operatorname{tr}(\mathbf{A}\mathbf{X})=\operatorname{tr}(\mathbf{A}\,\mathrm{d}\mathbf{X})$，再对照 $\operatorname{tr}(\mathbf{G}^\top\mathrm{d}\mathbf{X})$ 识别 $\mathbf{G}$）。
  \item （选读·克罗内克）用混合积律 \cref{eq:matderiv-kron-mixed} 证 $(\mathbf{A}\otimes\mathbf{B})^{-1}=\mathbf{A}^{-1}\otimes\mathbf{B}^{-1}$（提示：把 $(\mathbf{A}\otimes\mathbf{B})(\mathbf{A}^{-1}\otimes\mathbf{B}^{-1})$ 用 \cref{eq:matderiv-kron-mixed} 合并）；再由 \cref{eq:matderiv-vec-AXB} 写出 $\partial\operatorname{vec}(\mathbf{A}\mathbf{X}\mathbf{B})/\partial\operatorname{vec}(\mathbf{X})^\top$。
  \item （选读·舒尔补）对 $\mathbf{M}=\big[\begin{smallmatrix}\mathbf{A}&\mathbf{B}\\\mathbf{C}&\mathbf{D}\end{smallmatrix}\big]$ 用 \cref{prop:matderiv-block-inv} 验证 $\det\mathbf{M}=\det\mathbf{D}\det(\mathbf{M}/\mathbf{D})$，并由 \cref{eq:matderiv-block-inv} 的左上块读出 SMW 引理 $(\mathbf{A}-\mathbf{B}\mathbf{D}^{-1}\mathbf{C})^{-1}$ 的展开。
  \item （选读·Van Loan）取标量随机游走 $\dot x=w$（即 $\mathbf{A}=0,\mathbf{L}=1$，PSD $=q$）。先由 \cref{eq:matderiv-PhiQ-def} 手算 $\Phi=1,\ Q_k=q\,\Delta t$；再用 \cref{thm:matderiv-vanloan} 对 $\boldsymbol{\Xi}=\big[\begin{smallmatrix}0&q\\0&0\end{smallmatrix}\big]$ 取指数，验证读出 $\boldsymbol{\Upsilon}_{12}=q\,\Delta t,\ \boldsymbol{\Upsilon}_{22}=1$，得同样的 $\Phi,Q_k$——并对照 \cref{ch:eskf}/IMU 章的 $\sigma_d=\sigma/\sqrt{\Delta t}$ 白噪声换算。
  \item （选读·FIM）设标量参数 $\theta>0$ 重参数化为 $\phi=\ln\theta$（对数尺度），$J=\partial\theta/\partial\phi=\theta$。用 \cref{thm:matderiv-fim-reparam} 写出 $\mathcal{I}_\phi=\theta^2\mathcal{I}_\theta$，并由 \cref{cor:matderiv-crlb-reparam} 给 $\hat\phi$ 的 CRLB；解释为何“对数尺度上的相对误差”常更适合作不确定度度量。
\end{enumerate}
\end{practice}
